\chapter{Pupila circular uniforme: resultados analíticos}
\label{apx:focal_forms}

Este apéndice reúne los resultados explícitos para el campo en el plano focal de un sistema estigmático refractivo cuando el campo incidente es constante sobre una pupila circular de radio $a$. Se presentan las tres representaciones de polarización estudiadas en el capítulo~\ref{ch:focal}: circular, cilíndrica y cartesiana. En cada caso se parte de la forma general en transformadas de Hankel y se sustituyen las aproximaciones paraxiales de los coeficientes de Fresnel.

% -----------------------------------------------------------------------
\section{Representación circular}
% -----------------------------------------------------------------------


El punto de partida es el resultado exacto \eqref{eq:E_L_E_R_Hankel} del capítulo~\ref{ch:focal},

\begin{equation}
\begin{aligned}
E'_L(\vb r,f)
&=
\frac{\pi n}{i\lambda f z_0}
e^{inkf}
e^{i\frac{nk}{2f}r^2}
\Big[
\mathscr H_0\{t_+\mathscr E_L\}(q)
-
e^{-2i\varphi}
\mathscr H_2\{t_-\mathscr E_R\}(q)
\Big],
\\
E'_R(\vb r,f)
&=
\frac{\pi n}{i\lambda f z_0}
e^{inkf}
e^{i\frac{nk}{2f}r^2}
\Big[
-e^{2i\varphi} \mathscr H_2\{t_-\mathscr E_L\}(q)
+
\mathscr H_0\{t_+\mathscr E_R\}(q)
\Big],
\end{aligned}
\label{eq:E_L_E_R_Hankel_recall}
\end{equation}

donde $\mathscr H_m$ es la transformada de Hankel de orden $m$ definida en
\eqref{eq:def_Hankel_Transform_m_order} y $t_\pm=t_p\pm t_s$ según
\eqref{eq:t_plus_t_minus_def}.

El término transformado por $\mathscr{H}_0$ es diagonal y no mezcla las helicidades: puede identificarse con el correspondiente a la teoría escalar de la difracción, salvo por la modulación por $t_+$; si esta se toma constante, se obtiene exactamente el resultado escalar.

Los términos que $\mathscr{H}_2$ transforma acoplan $\ket{L}\leftrightarrow\ket{R}$ con fase azimutal $e^{\mp 2i\varphi}$ y, dado que $J_2(0)=0$, su contribución es estrictamente fuera de eje, por lo que dicha contribución ensancha el tamaño de la mancha focal de manera general y da lugar a un cambio en la polarización.

Las series paraxiales, derivadas en el apéndice~\ref{apx:ovoids_series} (véanse \eqref{eq:tp_plus_ts_series_fourth_order} y \eqref{eq:tp_minus_ts_series_fourth_order}), son, al orden cuadrático,

\begin{equation}
t_+(\rho')
\simeq
\frac{4n_0}{n_0+n_i}
+
(n_0+n_i)\,\xi\,\rho'^2,
\qquad
t_-(\rho')
\simeq
(n_0-n_i)\,\xi\,\rho'^2,
\label{eq:t_plus_t_minus_paraxial}
\end{equation}

con $\xi$ definida en \eqref{eq:xi_definition_fresnel_series},

\begin{equation}
\xi
=
\frac{
n_0(n_0-n_i)(z_0-z_i)^2
}{
z_0^2 z_i^2 (n_0+n_i)^3
}.
\label{eq:xi_focal_forms}
\end{equation}

Nótese que $t_+(0)=4n_0/(n_0+n_i)\neq 0$ mientras que $t_-(0)=0$: el canal de mezcla se enciende cuadráticamente al alejarse del eje.

Sustituyendo \eqref{eq:t_plus_t_minus_paraxial} en \eqref{eq:E_L_E_R_Hankel_recall}, se obtiene

\begin{equation}
\begin{aligned}
E'_L
\simeq{}&
\frac{\pi n}{i\lambda f z_0}
e^{inkf}
e^{i\frac{nk}{2f}r^2}
\left[
\frac{4n_0}{n_0+n_i}\mathscr H_0\{\mathscr E_L\}
+
(n_0+n_i)\xi\,\mathscr H_0\{r'^2\mathscr E_L\}
\right.
\\
&\qquad\qquad
\left.
-\,
(n_0-n_i)\xi\,e^{-2i\varphi}\mathscr H_2\{r'^2\mathscr E_R\}
\right],
\\[1em]
E'_R
\simeq{}&
\frac{\pi n}{i\lambda f z_0}
e^{inkf}
e^{i\frac{nk}{2f}r^2}
\left[
-(n_0-n_i)\xi\,e^{2i\varphi}\mathscr H_2\{r'^2\mathscr E_L\}
+
\frac{4n_0}{n_0+n_i}\mathscr H_0\{\mathscr E_R\}
\right.
\\
&\qquad\qquad
\left.
+\,
(n_0+n_i)\xi\,\mathscr H_0\{r'^2\mathscr E_R\}
\right].
\end{aligned}
\label{eq:E_L_E_R_paraxial}
\end{equation}


Si además el campo incidente es constante sobre la pupila circular de radio $a$, las funciones $\mathscr E_L$ y $\mathscr E_R$ son

\begin{equation}
  \mathscr E_L(r')=E_{L,0}A_a(r'),
  \qquad
  \mathscr E_R(r')=E_{R,0}A_a(r'),
\end{equation}

con $A_a$ la función pupila definida en \eqref{eq:def_pupil_function}. Las transformadas de Hankel necesarias se obtienen con facilidad por el uso de las bien conocidas identidades de las funciones de Bessel \cite{arfken2013}:

\begin{equation}
  \mathscr H_0\{A_a\}(q)
  =
  \frac{a}{q}J_1(qa),
  \label{eq:H0_chiR}
\end{equation}

\begin{equation}
  \mathscr H_0\{r'^2A_a\}(q)
  =
  \frac{a^3}{q}J_1(qa)
  +
  \frac{2a^2}{q^2}J_0(qa)
  -
  \frac{4a}{q^3}J_1(qa),
  \label{eq:H0_rho2_chiR}
\end{equation}

\begin{equation}
  \mathscr H_2\{r'^2A_a\}(q)
  =
  \frac{a^3}{q}J_3(qa).
  \label{eq:H2_rho2_chiR}
\end{equation}


Sustituyendo directamente en \eqref{eq:E_L_E_R_paraxial}, tenemos que

\begin{equation}
  \begin{aligned}
    E'_L(r,\varphi,f)
    &\simeq
    \frac{\pi n}{i\lambda f z_0}
    e^{inkf}
    e^{i\frac{nk}{2f}r^2}
    \left[
      \mathcal A(q)\,E_{L,0}
      -
      e^{-2i\varphi}\,\mathcal B(q)\,E_{R,0}
    \right],
    \\[0.5em]
    E'_R(r,\varphi,f)
    &\simeq
    \frac{\pi n}{i\lambda f z_0}
    e^{inkf}
    e^{i\frac{nk}{2f}r^2}
    \left[
      -
      e^{2i\varphi}\,\mathcal B(q)\,E_{L,0}
      +
      \mathcal A(q)\,E_{R,0}
    \right],
  \end{aligned}
  \label{eq:EL_ER_constant_E0}
\end{equation}

donde las funciones radiales $\mathcal A$ y $\mathcal B$ son

\begin{equation}
  \begin{aligned}
    \mathcal A(q)
    &=
    \frac{4n_0}{n_0+n_i}\,
    \frac{a}{q}J_1(qa)
    +
    (n_0+n_i)\,\xi
    \left(
      \frac{a^3}{q}J_1(qa)
      +
      \frac{2a^2}{q^2}J_0(qa)
      -
      \frac{4a}{q^3}J_1(qa)
    \right),
    \\[0.5em]
    \mathcal B(q)
    &=
    (n_0-n_i)\,\xi\,
    \frac{a^3}{q}J_3(qa).
  \end{aligned}
  \label{eq:A_B_circular_def}
\end{equation}

En forma matricial, con la matriz de mezcla $\mathbb V(\varphi)$ definida en \eqref{eq:def_V_mixing},

\begin{equation}
  \begin{pmatrix}
    E'_L(r,\varphi,f) \\
    E'_R(r,\varphi,f)
  \end{pmatrix}
  \simeq
  \frac{\pi n}{i\lambda f z_0}
  e^{inkf}
  e^{i\frac{nk}{2f}r^2}
  \left[
    \mathcal A(q)\,\mathbb I_\perp
    -
    \mathcal B(q)\,\mathbb V(\varphi)
  \right]
  \begin{pmatrix}
    E_{L,0} \\
    E_{R,0}
  \end{pmatrix}.
  \label{eq:EL_ER_constant_E0_matrix}
\end{equation}

Agrupando por funciones de Bessel, las funciones radiales se reescriben como

\begin{equation}
  \begin{aligned}
    \mathcal A(q)
    &=
    \frac{2(n_0+n_i)\xi a^2}{q^2}\,J_0(qa)
    +
    \left[
      \frac{4n_0}{n_0+n_i}\,\frac{a}{q}
      +
      (n_0+n_i)\xi
      \left(
        \frac{a^3}{q}
        -
        \frac{4a}{q^3}
      \right)
    \right] J_1(qa),
    \\[0.5em]
    \mathcal B(q)
    &=
    (n_0-n_i)\xi\,\frac{a^3}{q}\,J_3(qa).
  \end{aligned}
  \label{eq:EL_ER_constant_E0_bessel_series}
\end{equation}

En términos de los parámetros físicos del sistema, las combinaciones que pesan cada canal son

\begin{equation}
  (n_0+n_i)\,\xi
  =
  \frac{
    n_0(n_0-n_i)(z_0-z_i)^2
  }{
    z_0^2 z_i^2 (n_0+n_i)^2
  },
  \qquad
  (n_0-n_i)\,\xi
  =
  \frac{
    n_0(n_0-n_i)^2(z_0-z_i)^2
  }{
    z_0^2 z_i^2 (n_0+n_i)^3
  }.
  \label{eq:xi_combinations_physical}
\end{equation}

La estructura de \eqref{eq:EL_ER_constant_E0_matrix} separa los dos canales. La función $\mathcal A$, dominada por el término $\tfrac{4n_0}{n_0+n_i}\tfrac{a}{q}J_1(qa)$, es el patrón de Airy de la teoría escalar modulado por $t_+(0)$, con una corrección cuadrática de peso $(n_0+n_i)\xi$; como $\mathcal A(0)\neq 0$, este canal controla el pico central. La función $\mathcal B\propto J_3(qa)/q$ se anula en $q=0$, de modo que la conversión de helicidad —acompañada de las fases azimutales $e^{\mp 2i\varphi}$— es estrictamente fuera de eje y de peso $(n_0-n_i)\xi$: esta es la manifestación de la aberración vectorial en la representación circular.

% -----------------------------------------------------------------------
\section{Representación cilíndrica}
% -----------------------------------------------------------------------


En esta representación no hay acoplamiento entre las componentes radial y azimutal: la diagonalidad de \(\mathbb T\) se conserva en la reducción axialmente simétrica, pero la propagación difractiva convierte ambas componentes en transformadas de Hankel de orden uno.

El hecho de que aparezca \(J_1\) implica que el campo focal no posee contribución central para un campo incidente regular y axialmente simétrico en la base cilíndrica. En efecto, cerca de \(q=0\), las transformadas de Hankel de orden uno se anulan linealmente en \(q\). Por tanto,

\begin{equation}
  \vb E'_\perp(\vb 0,f)=0,
\end{equation}

siempre que las integrales radiales correspondientes sean finitas.


Para obtener resultados en forma cerrada, especializamos el campo incidente a una pupila circular uniforme y aproximamos los coeficientes de Fresnel por su expansión paraxial. Para esta última, conectamos con las series del apéndice~\ref{apx:ovoids_series} (véanse \eqref{eq:fresnel_ts_series_fourth_order} y \eqref{eq:fresnel_tp_series_fourth_order}). Con $\xi$ definida en \eqref{eq:xi_focal_forms}, al orden cuadrático

\begin{equation}
  t_p(\rho')
  \simeq
  \frac{2n_0}{n_0+n_i}
  +
  n_0\xi\,\rho'^2,
  \label{eq:tp_paraxial_cylindrical}
\end{equation}

\begin{equation}
  t_s(\rho')
  \simeq
  \frac{2n_0}{n_0+n_i}
  +
  n_i\xi\,\rho'^2.
  \label{eq:ts_paraxial_cylindrical}
\end{equation}

Sustituyendo en \eqref{eq:E_field_cylindrical_Hankel},

\begin{equation}
\begin{aligned}
E'_r
&\simeq
-\frac{2\pi n}{\lambda f z_0}
e^{inkf}
e^{i\frac{nk}{2f}r^2}
\left[
\frac{2n_0}{n_0+n_i}\mathscr H_1\{\mathscr E_r\}
+
n_0\xi
\mathscr H_1\{\rho'^2\mathscr E_r\}
\right],
\\
E'_\phi
&\simeq
-\frac{2\pi n}{\lambda f z_0}
e^{inkf}
e^{i\frac{nk}{2f}r^2}
\left[
\frac{2n_0}{n_0+n_i}\mathscr H_1\{\mathscr E_\phi\}
+
n_i\xi
\mathscr H_1\{\rho'^2\mathscr E_\phi\}
\right].
\end{aligned}
\label{eq:E_cylindrical_paraxial}
\end{equation}

Consideremos una pupila circular de radio \(a\), definida por

\begin{equation}
  A_a(\rho')
  =
  \begin{cases}
  1, & 0\leq \rho'\leq a,\\
  0, & \rho'>a.
  \end{cases}
\end{equation}

Tomamos

\begin{equation}
  \mathscr E_r(\rho')=E_{r,0}A_a(\rho'),
  \qquad
  \mathscr E_\phi(\rho')=E_{\phi,0}A_a(\rho').
\end{equation}

Las transformadas necesarias son

\begin{equation}
  \mathscr H_1\{A_a\}(q)
  =
  \int_0^a
  \rho'J_1(q\rho')\,d\rho',
\end{equation}

y

\begin{equation}
  \mathscr H_1\{\rho'^2A_a\}(q)
  =
  \int_0^a
  \rho'^3J_1(q\rho')\,d\rho'.
\end{equation}

Definimos

\begin{equation}
  \mathcal I(x)
  =
  \int_0^x J_0(u)\,du.
  \label{eq:Lommel_I_def_cylindrical}
\end{equation}

Con \(x=qa\), se obtiene

\begin{equation}
  L_1(q)
  \equiv
  \mathscr H_1\{A_a\}(q)
  =
  \frac{\mathcal I(qa)-qaJ_0(qa)}{q^2},
  \label{eq:L1_def_cylindrical}
\end{equation}

y

\begin{equation}
  M_1(q)
  \equiv
  \mathscr H_1\{\rho'^2A_a\}(q)
  =
  \frac{
  (qa)^3J_2(qa)
  +
  (qa)^2J_1(qa)
  +
  3qaJ_0(qa)
  -
  3\mathcal I(qa)
  }{q^4}.
  \label{eq:M1_def_cylindrical}
\end{equation}

Por tanto,

\begin{equation}
\begin{aligned}
\vb E'_\perp(\vb r,f)
\simeq
-\frac{2\pi n}{\lambda f z_0}
e^{inkf}
e^{i\frac{nk}{2f}r^2}
\Bigg\{
&
E_{r,0}
\left[
\frac{2n_0}{n_0+n_i}L_1(q)
+
n_0\xi\,M_1(q)
\right]\uvec r(\varphi)
\\
+
&
E_{\phi,0}
\left[
\frac{2n_0}{n_0+n_i}L_1(q)
+
n_i\xi\,M_1(q)
\right]\uvec\phi(\varphi)
\Bigg\}.
\end{aligned}
\label{eq:E_cylindrical_homogeneous}
\end{equation}

Los límites en el eje se obtienen directamente de las series de Bessel:

\begin{equation}
  L_1(q)
  =
  \frac{qa^3}{6}
  +
  O(q^3),
  \qquad
  M_1(q)
  =
  \frac{qa^5}{10}
  +
  O(q^3).
  \label{eq:L1_M1_small_q}
\end{equation}

De este modo, cada componente del campo transversal es de orden \(O(q)\) cerca del origen del plano focal. La contribución transversal \(I_\perp\), por consiguiente, se anula como \(O(q^2)\).


Como \(\uvec r(\varphi)\) y \(\uvec\phi(\varphi)\) son ortogonales punto a punto,

\begin{equation}
  \uvec r(\varphi)\cdot\uvec\phi(\varphi)=0,
\end{equation}

la intensidad transversal se obtiene sin término de interferencia entre las componentes radial y azimutal:

\begin{equation}
  I_\perp(\vb r,f)
  =
  \left(
  \frac{2\pi n}{\lambda f z_0}
  \right)^2
  \left[
  \left|
  \mathscr H_1\{t_p\mathscr E_r\}(q)
  \right|^2
  +
  \left|
  \mathscr H_1\{t_s\mathscr E_\phi\}(q)
  \right|^2
  \right].
\label{eq:I_perp_cylindrical}
\end{equation}

Sin embargo, la intensidad total debe incluir la componente longitudinal. Usando

\begin{equation}
E'_z(\vb r,f)=-\tan\theta_A'(r)\,E'_r(\vb r,f),
\qquad
\theta_A'(r)=\frac{r}{z_i-z(r)},
\label{eq:Ez_from_Er_cylindrical}
\end{equation}

se obtiene

\begin{equation}
I_z(\vb r,f)=|E'_z|^2=\tan^2\theta_A'(r)\,|E'_r|^2,
\label{eq:Iz_cylindrical}
\end{equation}

y por tanto

\begin{equation}
I(\vb r,f)=I_\perp(\vb r,f)+I_z(\vb r,f)
=
\left(
\frac{2\pi n}{\lambda f z_0}
\right)^2
\left[
\left(1+\tan^2\theta_A'(r)\right)
\left|
\mathscr H_1\{t_p\mathscr E_r\}(q)
\right|^2
+
\left|
\mathscr H_1\{t_s\mathscr E_\phi\}(q)
\right|^2
\right].
\label{eq:I_total_cylindrical}
\end{equation}

Para el caso uniforme sobre la pupila,

\begin{equation}
\begin{aligned}
\mathcal G_r(q)
\equiv\;&
\frac{2n_0}{n_0+n_i}
\frac{\mathcal I(qa)-qaJ_0(qa)}{q^2}
\\[0.3em]
&+
\frac{
n_0^2(n_0-n_i)(z_0-z_i)^2
}{
z_0^2z_i^2(n_0+n_i)^3
}
\frac{
(qa)^3J_2(qa)
+
(qa)^2J_1(qa)
+
3qaJ_0(qa)
-
3\mathcal I(qa)
}{q^4},
\end{aligned}
\label{eq:Gr_cylindrical}
\end{equation}

\begin{equation}
\begin{aligned}
\mathcal G_\phi(q)
\equiv\;&
\frac{2n_0}{n_0+n_i}
\frac{\mathcal I(qa)-qaJ_0(qa)}{q^2}
\\[0.3em]
&+
\frac{
n_0n_i(n_0-n_i)(z_0-z_i)^2
}{
z_0^2z_i^2(n_0+n_i)^3
}
\frac{
(qa)^3J_2(qa)
+
(qa)^2J_1(qa)
+
3qaJ_0(qa)
-
3\mathcal I(qa)
}{q^4},
\end{aligned}
\label{eq:Gphi_cylindrical}
\end{equation}

\begin{equation}
\begin{aligned}
  I_\perp(\vb r,f)
  \simeq
  \left(
  \frac{2\pi n}{\lambda f z_0}
  \right)^2
  \Bigg[
  &
  |E_{r,0}|^2
  \left|
  \mathcal G_r(q)
  \right|^2
  \\
  +
  &
  |E_{\phi,0}|^2
  \left|
  \mathcal G_\phi(q)
  \right|^2
  \Bigg].
\end{aligned}
\label{eq:I_perp_cylindrical_homogeneous}
\end{equation}

incluyendo la componente longitudinal:

\begin{equation}
\begin{aligned}
I(\vb r,f)
&\simeq
\left(
\frac{2\pi n}{\lambda f z_0}
\right)^2
\Bigg[
\left(1+\tan^2\theta_A'(r)\right)
|E_{r,0}|^2
\left|
\mathcal G_r(q)
\right|^2
\\
&\qquad
+
|E_{\phi,0}|^2
\left|
\mathcal G_\phi(q)
\right|^2
\Bigg].
\end{aligned}
\label{eq:I_total_cylindrical_homogeneous}
\end{equation}

Este resultado expresa una diferencia estructural con la representación circular: en la base circular aparecen naturalmente los órdenes \(0\) y \(2\), asociados respectivamente al término diagonal y al acoplamiento \(\ket L\leftrightarrow\ket R\). En cambio, para campos axialmente simétricos escritos directamente en la base cilíndrica, la dependencia angular de los propios vectores \(\uvec r(\phi')\) y \(\uvec\phi(\phi')\) impone el orden \(1\). Por eso los haces radial y azimutalmente polarizados presentan un cero axial en la componente transversal del campo focal: la ausencia de campo en el centro es un resultado puramente vectorial que la teoría escalar no puede describir, pues emerge de la estructura geométrica de los vectores $\uvec r(\varphi')$ y $\uvec\phi(\varphi')$ sobre la pupila, no de una modulación de amplitud.


La ecuación \eqref{eq:E_cylindrical_paraxial} escrita en forma matricial es

\begin{equation}
\begin{pmatrix}
E'_r\\
E'_\phi
\end{pmatrix}
\simeq
-\frac{2\pi n}{\lambda f z_0}
e^{inkf}
e^{i\frac{nk}{2f}r^2}
\begin{pmatrix}
\dfrac{2n_0}{n_0+n_i}\mathscr H_1+n_0\xi\,\mathscr H_1[\rho'^2]
&
0
\\[0.8em]
0
&
\dfrac{2n_0}{n_0+n_i}\mathscr H_1+n_i\xi\,\mathscr H_1[\rho'^2]
\end{pmatrix}
\begin{pmatrix}
\mathscr E_r\\
\mathscr E_\phi
\end{pmatrix},
\label{eq:E_cylindrical_paraxial_matrix}
\end{equation}

donde $\mathscr H_1[\rho'^2]$ denota el operador que multiplica por $\rho'^2$ y aplica la transformada de Hankel de orden uno.


Sustituyendo las definiciones \eqref{eq:L1_def_cylindrical} y \eqref{eq:M1_def_cylindrical} en \eqref{eq:E_cylindrical_homogeneous}, las componentes resultan

\begin{equation}
\begin{aligned}
E'_r
&\simeq
-\frac{2\pi n}{\lambda f z_0}
e^{inkf}
e^{i\frac{nk}{2f}r^2}
E_{r,0}
\Bigg[
\frac{2n_0}{n_0+n_i}
\frac{\mathcal I(qa)-qaJ_0(qa)}{q^2}
\\
&\qquad\qquad
+
\frac{
n_0^2(n_0-n_i)(z_0-z_i)^2
}{
z_0^2z_i^2(n_0+n_i)^3
}
\frac{
(qa)^3J_2(qa)
+
(qa)^2J_1(qa)
+
3qaJ_0(qa)
-
3\mathcal I(qa)
}{q^4}
\Bigg],
\\[0.4em]
E'_\phi
&\simeq
-\frac{2\pi n}{\lambda f z_0}
e^{inkf}
e^{i\frac{nk}{2f}r^2}
E_{\phi,0}
\Bigg[
\frac{2n_0}{n_0+n_i}
\frac{\mathcal I(qa)-qaJ_0(qa)}{q^2}
\\
&\qquad\qquad
+
\frac{
n_0n_i(n_0-n_i)(z_0-z_i)^2
}{
z_0^2z_i^2(n_0+n_i)^3
}
\frac{
(qa)^3J_2(qa)
+
(qa)^2J_1(qa)
+
3qaJ_0(qa)
-
3\mathcal I(qa)
}{q^4}
\Bigg].
\end{aligned}
\label{eq:E_cylindrical_homogeneous_explicit}
\end{equation}

Agrupando por orden de función de Bessel, con $x = qa$:

\begin{equation}
\begin{aligned}
E'_r
&\simeq
-\frac{2\pi n}{\lambda f z_0}
e^{inkf}
e^{i\frac{nk}{2f}r^2}
E_{r,0}
\Bigg[
\left(
\frac{2n_0}{n_0+n_i}\frac{1}{q^2}
-
\frac{3n_0^2(n_0-n_i)(z_0-z_i)^2}{z_0^2z_i^2(n_0+n_i)^3}\frac{1}{q^4}
\right)\mathcal I(x)
\\
&\qquad
+
\left(
-\frac{2n_0}{n_0+n_i}\frac{x}{q^2}
+
\frac{3n_0^2(n_0-n_i)(z_0-z_i)^2}{z_0^2z_i^2(n_0+n_i)^3}\frac{x}{q^4}
\right)J_0(x)
+
\left(
\frac{n_0^2(n_0-n_i)(z_0-z_i)^2}{z_0^2z_i^2(n_0+n_i)^3}\frac{x^2}{q^4}
\right)J_1(x)
\\
&\qquad
+
\left(
\frac{n_0^2(n_0-n_i)(z_0-z_i)^2}{z_0^2z_i^2(n_0+n_i)^3}\frac{x^3}{q^4}
\right)J_2(x)
\Bigg],
\\[0.6em]
E'_\phi
&\simeq
-\frac{2\pi n}{\lambda f z_0}
e^{inkf}
e^{i\frac{nk}{2f}r^2}
E_{\phi,0}
\Bigg[
\left(
\frac{2n_0}{n_0+n_i}\frac{1}{q^2}
-
\frac{3n_0n_i(n_0-n_i)(z_0-z_i)^2}{z_0^2z_i^2(n_0+n_i)^3}\frac{1}{q^4}
\right)\mathcal I(x)
\\
&\qquad
+
\left(
-\frac{2n_0}{n_0+n_i}\frac{x}{q^2}
+
\frac{3n_0n_i(n_0-n_i)(z_0-z_i)^2}{z_0^2z_i^2(n_0+n_i)^3}\frac{x}{q^4}
\right)J_0(x)
+
\left(
\frac{n_0n_i(n_0-n_i)(z_0-z_i)^2}{z_0^2z_i^2(n_0+n_i)^3}\frac{x^2}{q^4}
\right)J_1(x)
\\
&\qquad
+
\left(
\frac{n_0n_i(n_0-n_i)(z_0-z_i)^2}{z_0^2z_i^2(n_0+n_i)^3}\frac{x^3}{q^4}
\right)J_2(x)
\Bigg].
\end{aligned}
\label{eq:E_cylindrical_homogeneous_grouped_bessel}
\end{equation}


El resultado \eqref{eq:E_cylindrical_homogeneous} escrito en forma matricial es

\begin{equation}
\begin{pmatrix}
E'_r\\
E'_\phi
\end{pmatrix}
\simeq
-\frac{2\pi n}{\lambda f z_0}
e^{inkf}
e^{i\frac{nk}{2f}r^2}
\begin{pmatrix}
\dfrac{2n_0}{n_0+n_i}L_1(q)+n_0\xi\,M_1(q) & 0
\\[0.8em]
0 & \dfrac{2n_0}{n_0+n_i}L_1(q)+n_i\xi\,M_1(q)
\end{pmatrix}
\begin{pmatrix}
E_{r,0}\\
E_{\phi,0}
\end{pmatrix}.
\label{eq:E_cylindrical_homogeneous_matrix}
\end{equation}

% -----------------------------------------------------------------------
\section{Representación cartesiana}
% -----------------------------------------------------------------------


Para obtener resultados en forma cerrada, especializamos el campo incidente a una pupila circular homogénea y empleamos la expansión paraxial de las combinaciones $t_+$ y $t_-$ derivada en el apéndice~\ref{apx:ovoids_series}. Como en \eqref{eq:t_plus_t_minus_paraxial},

\begin{equation}
t_+(\rho')
\simeq
\frac{4n_0}{n_0+n_i}
+
(n_0+n_i)\,\xi\,\rho'^2,
\qquad
t_-(\rho')
\simeq
(n_0-n_i)\,\xi\,\rho'^2,
\label{eq:t_plus_t_minus_paraxial_cartesian}
\end{equation}

con $\xi$ dada por \eqref{eq:xi_focal_forms}. Sustituyendo en \eqref{eq:Ex_Ey_Hankel_cartesian_matrix_separated}, la expresión compacta es

\begin{equation}
\begin{aligned}
\mathbf E'_\perp(\mathbf r,f)
\simeq{}&
\frac{\pi n e^{inkf}}{i\lambda f z_0}
e^{i\frac{nk}{2f}r^2}
\left[
\frac{4n_0}{n_0+n_i}
\mathscr H_0\{\mathscr E_\perp\}(q)
+
(n_0+n_i)\xi\,
\mathscr H_0\{\rho'^2\mathscr E_\perp\}(q)
\right.
\\
&\qquad\qquad
\left.
-\,
(n_0-n_i)\xi\,
\mathbb W(\varphi)\,
\mathscr H_2\{\rho'^2\mathscr E_\perp\}(q)
\right].
\end{aligned}
\label{eq:E_field_cartesian_paraxial}
\end{equation}

Consideremos una pupila circular homogénea de radio $a$, para la cual

\begin{equation}
\mathscr E_\perp(\rho')
=
\mathscr E_{\perp,0}A_a(\rho'),
\qquad
\mathscr E_{\perp,0}
=
\mathscr E_{x,0}\,\hat{\mathbf x}
+
\mathscr E_{y,0}\,\hat{\mathbf y}.
\label{eq:cartesian_homogeneous_pupil_def}
\end{equation}

Usaremos las siguientes transformadas de Hankel:

\begin{equation}
\mathscr H_0\{A_a\}(q)
=
\frac{a}{q}J_1(qa),
\label{eq:H0_Aa_cartesian}
\end{equation}

\begin{equation}
\mathscr H_0\{\rho'^2A_a\}(q)
=
\frac{a^3}{q}J_1(qa)
+
\frac{2a^2}{q^2}J_0(qa)
-
\frac{4a}{q^3}J_1(qa),
\label{eq:H0_rho2_Aa_cartesian}
\end{equation}

\begin{equation}
\mathscr H_2\{\rho'^2A_a\}(q)
=
\frac{a^3}{q}J_3(qa).
\label{eq:H2_rho2_Aa_cartesian}
\end{equation}

Definimos

\begin{equation}
\begin{aligned}
\mathcal A(q)
&=
\frac{4n_0}{n_0+n_i}\,
\frac{a}{q}J_1(qa)
+
(n_0+n_i)\,\xi
\left[
\frac{a^3}{q}J_1(qa)
+
\frac{2a^2}{q^2}J_0(qa)
-
\frac{4a}{q^3}J_1(qa)
\right],
\\[0.5em]
\mathcal B(q)
&=
(n_0-n_i)\,\xi\,
\frac{a^3}{q}J_3(qa),
\end{aligned}
\label{eq:A_B_cartesian_def}
\end{equation}

que coinciden con las funciones radiales \eqref{eq:A_B_circular_def} de la representación circular: los dos canales radiales son los mismos en ambas bases y solo cambia la matriz de mezcla, $\mathbb V(\varphi)$ en la circular y $\mathbb W(\varphi)$ en la cartesiana.

Entonces el campo en el plano focal adopta la forma compacta

\begin{equation}
\mathbf E'_\perp(\mathbf r,f)
\simeq
\frac{\pi n e^{inkf}}{i\lambda f z_0}
e^{i\frac{nk}{2f}r^2}
\left[
\mathcal A(q)\mathbb I_\perp
-
\mathcal B(q)\mathbb W(\varphi)
\right]
\mathscr E_{\perp,0}.
\label{eq:E_field_cartesian_homogeneous}
\end{equation}

Por componentes,

\begin{equation}
\begin{aligned}
E'_x(\mathbf r,f)
&\simeq
\frac{\pi n e^{inkf}}{i\lambda f z_0}
e^{i\frac{nk}{2f}r^2}
\left[
\left(
\mathcal A(q)-\mathcal B(q)\cos2\varphi
\right)\mathscr E_{x,0}
-
\mathcal B(q)\sin2\varphi\,\mathscr E_{y,0}
\right],
\\[1em]
E'_y(\mathbf r,f)
&\simeq
\frac{\pi n e^{inkf}}{i\lambda f z_0}
e^{i\frac{nk}{2f}r^2}
\left[
-
\mathcal B(q)\sin2\varphi\,\mathscr E_{x,0}
+
\left(
\mathcal A(q)+\mathcal B(q)\cos2\varphi
\right)\mathscr E_{y,0}
\right].
\end{aligned}
\label{eq:Ex_Ey_cartesian_homogeneous}
\end{equation}


Para una polarización incidente lineal en $\hat{\mathbf x}$,

\begin{equation}
\mathscr E_{\perp,0}
=
\begin{pmatrix}
1\\
0
\end{pmatrix},
\end{equation}

se obtiene

\begin{equation}
\begin{aligned}
E'_x
&\simeq
\frac{\pi n e^{inkf}}{i\lambda f z_0}
e^{i\frac{nk}{2f}r^2}
\left[
\mathcal A(q)
-
\mathcal B(q)\cos2\varphi
\right],
\\
E'_y
&\simeq
-
\frac{\pi n e^{inkf}}{i\lambda f z_0}
e^{i\frac{nk}{2f}r^2}
\mathcal B(q)\sin2\varphi.
\end{aligned}
\label{eq:cartesian_case_10}
\end{equation}

Para una polarización incidente lineal en $\hat{\mathbf y}$,

\begin{equation}
\mathscr E_{\perp,0}
=
\begin{pmatrix}
0\\
1
\end{pmatrix},
\end{equation}

se obtiene

\begin{equation}
\begin{aligned}
E'_x
&\simeq
-
\frac{\pi n e^{inkf}}{i\lambda f z_0}
e^{i\frac{nk}{2f}r^2}
\mathcal B(q)\sin2\varphi,
\\
E'_y
&\simeq
\frac{\pi n e^{inkf}}{i\lambda f z_0}
e^{i\frac{nk}{2f}r^2}
\left[
\mathcal A(q)
+
\mathcal B(q)\cos2\varphi
\right].
\end{aligned}
\label{eq:cartesian_case_01}
\end{equation}

Para una polarización incidente diagonal,

\begin{equation}
\mathscr E_{\perp,0}
=
\begin{pmatrix}
1\\
1
\end{pmatrix},
\end{equation}

el campo focal queda

\begin{equation}
\begin{aligned}
E'_x
&\simeq
\frac{\pi n e^{inkf}}{i\lambda f z_0}
e^{i\frac{nk}{2f}r^2}
\left[
\mathcal A(q)
-
\mathcal B(q)
\left(
\cos2\varphi+\sin2\varphi
\right)
\right],
\\
E'_y
&\simeq
\frac{\pi n e^{inkf}}{i\lambda f z_0}
e^{i\frac{nk}{2f}r^2}
\left[
\mathcal A(q)
+
\mathcal B(q)
\left(
\cos2\varphi-\sin2\varphi
\right)
\right].
\end{aligned}
\label{eq:cartesian_case_11}
\end{equation}

Finalmente, para una polarización incidente

\begin{equation}
\mathscr E_{\perp,0}
=
\begin{pmatrix}
1\\
i
\end{pmatrix},
\end{equation}

se obtiene

\begin{equation}
\begin{aligned}
E'_x
&\simeq
\frac{\pi n e^{inkf}}{i\lambda f z_0}
e^{i\frac{nk}{2f}r^2}
\left[
\mathcal A(q)
-
\mathcal B(q)
\left(
\cos2\varphi+i\sin2\varphi
\right)
\right],
\\
E'_y
&\simeq
\frac{\pi n e^{inkf}}{i\lambda f z_0}
e^{i\frac{nk}{2f}r^2}
\left[
i\mathcal A(q)
-
\mathcal B(q)
\left(
\sin2\varphi-i\cos2\varphi
\right)
\right].
\end{aligned}
\label{eq:cartesian_case_1i}
\end{equation}

Usando

\begin{equation}
  \cos2\varphi+i\sin2\varphi
  =
  e^{2i\varphi},
\end{equation}

la primera componente del caso $(1,i)$ puede escribirse como

\begin{equation}
E'_x
\simeq
\frac{\pi n e^{inkf}}{i\lambda f z_0}
e^{i\frac{nk}{2f}r^2}
\left[
\mathcal A(q)
-
\mathcal B(q)e^{2i\varphi}
\right].
\end{equation}

Esto muestra explícitamente cómo, en el caso circular escrito desde la base
cartesiana, reaparece la fase azimutal de orden dos.


La estructura obtenida en la base cartesiana es

\begin{equation}
  \mathscr H_0\,\mathbb I_\perp
  -
  \mathscr H_2\,\mathbb W(\varphi).
\end{equation}

El término proporcional a $\mathscr H_0$ es la contribución isotrópica:
preserva las componentes cartesianas incidentes y corresponde al canal de
orden angular cero. En ausencia de la modulación vectorial introducida por
$t_-$, este término reproduce la estructura central esperada de la teoría
escalar de la difracción, salvo por la modulación radial debida a los
coeficientes de Fresnel.

El término proporcional a $\mathscr H_2$ contiene la información vectorial de
la transmisión. Está gobernado por la diferencia

\begin{equation}
  t_-
  =
  t_p-t_s,
\end{equation}

y por el operador angular

\begin{equation}
  \mathbb W(\varphi)
  =
  \begin{pmatrix}
    \cos 2\varphi & \sin 2\varphi \\
    \sin 2\varphi & -\cos 2\varphi
  \end{pmatrix}.
\end{equation}

Por tanto, este término mezcla las componentes $x$ e $y$ del campo propagado
y produce una redistribución angular de orden dos en el plano focal.

La misma física que en la representación circular aparece como acoplamiento
helicoidal entre $\ket L$ y $\ket R$ mediante fases $e^{\pm2i\varphi}$, en la
representación cartesiana aparece como una anisotropía angular gobernada por
$\cos2\varphi$ y $\sin2\varphi$. En síntesis,

\begin{equation}
\text{representación circular:}
\qquad
\mathscr H_0\,\mathbb I_\perp
-
\mathscr H_2\,\mathbb V(\varphi),
\end{equation}

mientras que

\begin{equation}
\text{representación cartesiana:}
\qquad
\mathscr H_0\,\mathbb I_\perp
-
\mathscr H_2\,\mathbb W(\varphi).
\end{equation}

Ambas expresiones describen el mismo fenómeno vectorial: la parte escalar
conserva la polarización, mientras que la parte de orden angular dos introduce
mezcla polarizacional y estructura azimutal en el plano focal.


La expresión compacta \eqref{eq:E_field_cartesian_paraxial} escrita por componentes es

\begin{equation}
\begin{aligned}
E'_x(\mathbf r,f)
&\simeq
\frac{\pi n e^{inkf}}{i\lambda f z_0}
e^{i\frac{nk}{2f}r^2}
\Bigg[
\frac{4n_0}{n_0+n_i}
\mathscr H_0\{\mathscr E_x\}(q)
+
(n_0+n_i)\xi\,
\mathscr H_0\{\rho'^2\mathscr E_x\}(q)
\\
&\quad
-
(n_0-n_i)\xi
\left(
\cos2\varphi\,
\mathscr H_2\{\rho'^2\mathscr E_x\}(q)
+
\sin2\varphi\,
\mathscr H_2\{\rho'^2\mathscr E_y\}(q)
\right)
\Bigg],
\\[1em]
E'_y(\mathbf r,f)
&\simeq
\frac{\pi n e^{inkf}}{i\lambda f z_0}
e^{i\frac{nk}{2f}r^2}
\Bigg[
\frac{4n_0}{n_0+n_i}
\mathscr H_0\{\mathscr E_y\}(q)
+
(n_0+n_i)\xi\,
\mathscr H_0\{\rho'^2\mathscr E_y\}(q)
\\
&\quad
-
(n_0-n_i)\xi
\left(
\sin2\varphi\,
\mathscr H_2\{\rho'^2\mathscr E_x\}(q)
-
\cos2\varphi\,
\mathscr H_2\{\rho'^2\mathscr E_y\}(q)
\right)
\Bigg].
\end{aligned}
\label{eq:Ex_Ey_cartesian_paraxial}
\end{equation}


El resultado \eqref{eq:E_field_cartesian_homogeneous} escrito en forma matricial es

\begin{equation}
\begin{pmatrix}
E'_x\\
E'_y
\end{pmatrix}
\simeq
\frac{\pi n e^{inkf}}{i\lambda f z_0}
e^{i\frac{nk}{2f}r^2}
\begin{pmatrix}
\mathcal A(q)-\mathcal B(q)\cos2\varphi
&
-\mathcal B(q)\sin2\varphi
\\
-\mathcal B(q)\sin2\varphi
&
\mathcal A(q)+\mathcal B(q)\cos2\varphi
\end{pmatrix}
\begin{pmatrix}
\mathscr E_{x,0}\\
\mathscr E_{y,0}
\end{pmatrix}.
\label{eq:E_field_cartesian_homogeneous_matrix}
\end{equation}


Sustituyendo explícitamente \eqref{eq:A_B_cartesian_def} en \eqref{eq:Ex_Ey_cartesian_homogeneous},

\begin{equation}
\begin{aligned}
E'_x(\mathbf r,f)
&\simeq
\frac{\pi n e^{inkf}}{i\lambda f z_0}
e^{i\frac{nk}{2f}r^2}
\left[
\mathcal A(q)\,\mathscr E_{x,0}
-
\mathcal B(q)
\left(
\cos2\varphi\,\mathscr E_{x,0}
+
\sin2\varphi\,\mathscr E_{y,0}
\right)
\right],
\\[1em]
E'_y(\mathbf r,f)
&\simeq
\frac{\pi n e^{inkf}}{i\lambda f z_0}
e^{i\frac{nk}{2f}r^2}
\left[
\mathcal A(q)\,\mathscr E_{y,0}
-
\mathcal B(q)
\left(
\sin2\varphi\,\mathscr E_{x,0}
-
\cos2\varphi\,\mathscr E_{y,0}
\right)
\right].
\end{aligned}
\label{eq:Ex_Ey_cartesian_homogeneous_AB}
\end{equation}


O en forma de series de las funciones de Bessel: definiendo los factores angulares

\begin{equation}
\mathcal C(\varphi)
\equiv
\cos2\varphi\,\mathscr E_{x,0}+\sin2\varphi\,\mathscr E_{y,0},
\qquad
\mathcal S(\varphi)
\equiv
\sin2\varphi\,\mathscr E_{x,0}-\cos2\varphi\,\mathscr E_{y,0},
\label{eq:cartesian_angular_factors}
\end{equation}

y agrupando por orden, mediante \eqref{eq:EL_ER_constant_E0_bessel_series},

\begin{equation}
\begin{aligned}
E'_x(\mathbf r,f)
\simeq&\;
\frac{\pi n e^{inkf}}{i\lambda f z_0}
e^{i\frac{nk}{2f}r^2}
\Bigg\{
J_0(qa)\,
\frac{2(n_0+n_i)\xi a^2}{q^2}\,
\mathscr E_{x,0}
\\
&\qquad
+
J_1(qa)
\left[
\frac{4n_0}{n_0+n_i}\,\frac{a}{q}
+
(n_0+n_i)\xi
\left(
\frac{a^3}{q}
-
\frac{4a}{q^3}
\right)
\right]
\mathscr E_{x,0}
\\
&\qquad
-
J_3(qa)\,
(n_0-n_i)\xi\,\frac{a^3}{q}\,
\mathcal C(\varphi)
\Bigg\},
\\[1em]
E'_y(\mathbf r,f)
\simeq&\;
\frac{\pi n e^{inkf}}{i\lambda f z_0}
e^{i\frac{nk}{2f}r^2}
\Bigg\{
J_0(qa)\,
\frac{2(n_0+n_i)\xi a^2}{q^2}\,
\mathscr E_{y,0}
\\
&\qquad
+
J_1(qa)
\left[
\frac{4n_0}{n_0+n_i}\,\frac{a}{q}
+
(n_0+n_i)\xi
\left(
\frac{a^3}{q}
-
\frac{4a}{q^3}
\right)
\right]
\mathscr E_{y,0}
\\
&\qquad
-
J_3(qa)\,
(n_0-n_i)\xi\,\frac{a^3}{q}\,
\mathcal S(\varphi)
\Bigg\},
\end{aligned}
\label{eq:Ex_Ey_cartesian_bessel_series}
\end{equation}

donde las combinaciones $(n_0\pm n_i)\,\xi$ en términos de los parámetros
físicos del sistema están dadas en \eqref{eq:xi_combinations_physical}.

Las tres representaciones describen el mismo fenómeno físico: la transmisión
en la interfaz impone sobre la pupila una modulación no uniforme —el canal
$t_-$— que redistribuye energía fuera del eje respecto a la predicción
escalar. La elección de base determina únicamente cómo se manifiesta dicha
modulación: como mezcla de helicidades con fases $e^{\mp2i\varphi}$ en la
base circular, como acoplamiento entre componentes con anisotropía azimutal
de orden dos en la cartesiana, o como modulación radial pura en la
cilíndrica; el campo resultante es en todos los casos el mismo.
